% !PRINTED 580
\weberpage{580}

%<<CONT
follows, when we form the derivative with respect to $t$, which is evidently permissible:
\[
F'(t) \equiv e\,[P(t)]^{e-1}\,P'(t)\,\Phi(t) + [P(t)]^e\,\Phi'(t)\ (\mathrm{mod}\ p).
\]
If we set $t=\tau$ in this, then $P'(t)$ passes over into a form $P'(\tau)$ not divisible by $\mathfrak{p}$ (which we have already shown and used in the proof of \S.\,157, 4.), and we obtain
% [sic: Weber prints $F_1(\tau)$ here; the preceding factorization and following sentence use $\Phi(\tau)$.]
\weq{11}{F'(\tau) \equiv e\, P(\tau)^{e-1}\, P'(\tau)\, F_1(\tau)\ (\mathrm{mod}\ \mathfrak{p}^e).}
Now (by \S.\,157, 4.) $P(\tau)$ is divisible by $\mathfrak{p}$ but not by $\mathfrak{p}^2$, while $P'(\tau)$ and $\Phi(\tau)$ are not divisible by $\mathfrak{p}$, and so the formula (11) thus gives us the proof of the following theorem:

1. \gesperrt{If $\mathfrak{p}$ is an arbitrary prime ideal, $p$ the natural prime divisible by $\mathfrak{p}$, and $\mathfrak{p}^e$ the highest power of $\mathfrak{p}$ that divides $p$, then the fundamental form $F'(\tau)$ is in every case divisible by $\mathfrak{p}^{e-1}$; if, further, the exponent $e$ is not divisible by $p$, then $F'(\tau)$ is not divisible by $\mathfrak{p}^e$; but if $e$ is divisible by $p$, then $F'(\tau)$ is divisible by $\mathfrak{p}^e$ and perhaps by still higher powers of $\mathfrak{p}$.}

If $e$ is greater than 1, then by the foregoing $F'(\tau)$ is divisible by $\mathfrak{p}$, and consequently $N_a[F'(\tau)]$, and hence also the fundamental number $\Delta$, is divisible by $p$.

But if all prime ideals $\mathfrak{p}$ divide $p$ only in the first power, then $F'(\tau)$ is relatively prime to $p$. If one then decomposes $F'(\tau)$ into its prime factors, there occurs among them none whose norm is a power of $p$, and consequently $N_a[F'(\tau)]$ and $\Delta$ are likewise not divisible by $p$. From this there then follows the theorem:

2. \gesperrt{A natural prime $p$ is then and only then divisible in the field $\Omega$ by the square of a prime factor, when $p$ divides the fundamental number of $\Omega$.}

From these considerations we can draw still other important conclusions.

If we solve the system of linear equations (1) with respect to $\omega_1, \omega_2, \ldots, \omega_n$, we obtain expressions with the denominator $U$, which, as we have seen, is a unit. If one substitutes these expressions in
\weq{12}{\omega = x_1\omega_1 + x_2\omega_2 + \cdots + x_n\omega_n,}
%CONT>>
